\documentclass[11pt,a4paper]{article}

\usepackage{../shared/cathedral-preamble}
\usepackage{listings}
\lstset{basicstyle=\ttfamily\small, breaklines=true, frame=single, backgroundcolor=\color{gray!8}}

\title{%
  Formalizing Analytic Number Theory in Lean~4:\\
  Foundations, Practices, and Lessons\\
  from the Cathedral Project%
}
\author{
  Jason Robert Gochanour
}
\date{May 2026}

\begin{document}

\maketitle

\begin{abstract}
We describe the practical techniques and logical foundations for
formalizing analytic number theory in Lean~4 with Mathlib, distilled
from the Cathedral project: a $\sim$60{,}500-line formalization of the
Nyman--Beurling criterion for the Riemann Hypothesis.

The first half covers foundations: the type-theoretic setting (CIC
with extensions), the trust chain from silicon to theorem, the
precise logical strength of the formalization (provable in $\ZFC$),
and a constructivity analysis showing the converse direction is
partially constructive while the forward direction is inherently
classical.

The second half covers practice: Mathlib API patterns that worked
well (interval integrals, spectral theory), patterns requiring
workarounds (infinite sums, measure theory), 5 API gaps requiring
axioms, 12 essential lemmas, project structure advice, and common
tactic patterns. We document 3 deprecated-API traps and provide
a verification checklist for large-scale proof engineering.
\end{abstract}

\tableofcontents

% ================================================================
\section{Introduction}
\label{sec:intro}
% ================================================================

Every machine-verified proof rests on a stack of trust. This paper
examines that stack for the Cathedral, a Lean~4 formalization
establishing $\RH \iff d_N^2 \to 0$, from two perspectives:
the \emph{foundational} (what are we trusting?) and the
\emph{practical} (how do we build it?).

\begin{center}
\begin{tabular}{@{}cl@{}}
\toprule
\textbf{Layer} & \textbf{What Must Be Trusted} \\
\midrule
5 & Mathematical axioms (1 on crown path, $\sim$75 total) \\
4 & Mathlib library ($\sim$1.4M LOC) \\
3 & Lean kernel axioms (3: \lean{propext}, \lean{choice}, \lean{Quot.sound}) \\
2 & Lean type-checker ($\sim$5{,}000 LOC of C++ kernel) \\
1 & Compiler, OS, hardware \\
\bottomrule
\end{tabular}
\end{center}

The Cathedral project encountered the challenge of using complex
analysis, real analysis, combinatorics, and linear algebra
\emph{together} in a single proof, with consistent type coercions
and compatible assumptions. Over 45 days it produced 227 active
Lean files (with 103 archived). This paper distills the lessons.

% ================================================================
\section{The Type-Theoretic Foundation}
\label{sec:foundation}
% ================================================================

\subsection{CIC with Extensions}

Lean~4 is based on the Calculus of Inductive Constructions (CIC)
with dependent function types, inductive types, a universe hierarchy,
proof irrelevance, and large elimination. Three kernel axioms extend
the core:

\begin{enumerate}
  \item \textbf{\lean{propext}}: $(P \iff Q) \to P = Q$. Essential
    for set-theoretic reasoning.
  \item \textbf{\lean{Classical.choice}}: For any nonempty type $A$,
    there exists $a : A$. Implies excluded middle.
  \item \textbf{\lean{Quot.sound}}: If $a \sim b$ then $[a] = [b]$
    in the quotient type.
\end{enumerate}

The system $\CIC + \lean{propext} + \lean{choice} + \lean{Quot.sound}$
is equiconsistent with $\ZFC$.

\subsection{The Four Axioms of the Crown}

The crown theorem reports exactly 4 axioms via \lean{\#print axioms}
(3 kernel + 1 mathematical):

\begin{center}
\begin{tabular}{@{}clcc@{}}
\toprule
\textbf{\#} & \textbf{Axiom} & \textbf{Origin} & \textbf{Tier} \\
\midrule
1 & \lean{propext} & Lean kernel & --- \\
2 & \lean{Classical.choice} & Lean kernel & --- \\
3 & \lean{Quot.sound} & Lean kernel & --- \\
4 & \lean{baez\_duarte\_forward} & Cathedral & BD (IMRN 2003) \\
\bottomrule
\end{tabular}
\end{center}

The sole mathematical axiom is the B\'aez-Duarte forward direction
(IMRN 2003, no.~36), provable in $\ZFC$ by standard methods.
It is a formalization gap, not a foundational gap.

\subsection{Constructivity Analysis}

The converse ($d_N^2 \to 0 \implies \RH$) is partially constructive:
the bound $d_N^2 \geq \delta > 0$ is computed explicitly given
$\rho$; only the existential over $\rho$ in $\neg\RH$ is
non-constructive. The forward direction is inherently classical,
requiring the Montgomery--Vaughan MVT and Perron's formula.

\subsection{The de Bruijn Criterion}

\begin{center}
\begin{tabular}{@{}lr@{}}
\toprule
\textbf{Component} & \textbf{Size (LOC)} \\
\midrule
Lean kernel (type-checker) & $\sim$5{,}000 \\
Lean elaborator + tactics & $\sim$100{,}000 \\
Mathlib library & $\sim$1{,}400{,}000 \\
Cathedral & $\sim$60{,}500 \\
\bottomrule
\end{tabular}
\end{center}

Only the kernel (5{,}000 LOC) is in the trusted base.

% ================================================================
\section{The Axiom-Theorem Gradient}
\label{sec:gradient}
% ================================================================

The Cathedral introduces a foundational concept: the
\emph{axiom-theorem gradient}.

\begin{enumerate}
  \item \textbf{Kernel axioms}: Trusted by design. (3)
  \item \textbf{Crown axioms}: On the critical path. (1)
  \item \textbf{Off-path axioms}: Don't affect the crown. ($\sim$74)
  \item \textbf{Theorems}: Machine-verified from Mathlib.
  \item \textbf{Archived axioms}: Previously active, now proved. (11)
\end{enumerate}

This gradient makes progress visible. Each axiom elimination moves
a statement from category~2/3 to category~4. The reduction from
56 crown axioms to 1 is a deleveraging of the proof's balance sheet.

% ================================================================
\section{Mathlib APIs That Worked Well}
\label{sec:good-apis}
% ================================================================

\subsection{Interval Integrals}

Mathlib's \lean{MeasureTheory.integral} and \lean{intervalIntegral}
work well for $L^2$ norms on $(0,1)$. Key lemma:
\lean{integral\_nonneg} with \lean{sq\_nonneg} for $d_N^2 \geq 0$.

\subsection{Matrix API}

Mathlib's matrix library is excellent for Gram matrices.
\lean{Matrix.PosDef.det\_pos} connects positive definiteness to
invertibility. \lean{Matrix.nonsing\_inv} provides the inverse.

\subsection{Spectral Theory}

For eigenvalue bounds, use the Rayleigh quotient approach via
\lean{Matrix.IsHermitian} (which handles the real case via
\lean{starRingEnd}).

\subsection{Convergence}

When possible, avoid \lean{Filter.Tendsto} and \lean{Filter.atTop}.
Algebraic bounds with explicit inequalities are easier to formalize.

% ================================================================
\section{Mathlib APIs That Required Workarounds}
\label{sec:workarounds}
% ================================================================

\subsection{Infinite Series}

\lean{tsum} requires \lean{Summable} hypotheses difficult to
establish for Dirichlet series. We used truncated \lean{Finset.sum}
with explicit error bounds.

\subsection{Fractional Part}

\lean{Int.fract} exists but its API is sparse for analytic number
theory (integrability, measurability). Expect to prove many basic
properties yourself.

\subsection{Measure Theory Coercions}

Moving between Bochner and Henstock--Kurzweil integrals requires
care via \lean{intervalIntegrable\_iff\_integrable\_Ioc}.

% ================================================================
\section{Five Mathlib Gaps Requiring Axioms}
\label{sec:gaps}
% ================================================================

\begin{center}
\begin{tabular}{@{}lp{5cm}l@{}}
\toprule
\textbf{Gap} & \textbf{What's Missing} & \textbf{Difficulty} \\
\midrule
Analytic continuation & Mellin integral extension & Medium \\
Contour integrals & Vertical line integration & Hard \\
Mertens function & $M(x) = \sum_{n \leq x}\mu(n)$ & Medium \\
Montgomery--Vaughan & Dirichlet polynomial MVT & Hard \\
Dedekind sums & Off-diagonal Vasyunin & Medium \\
\bottomrule
\end{tabular}
\end{center}

These five additions would eliminate 80\% of the Cathedral's axioms.

% ================================================================
\section{Twelve Essential Mathlib Lemmas}
\label{sec:lemmas}
% ================================================================

\begin{center}
\small
\begin{tabular}{@{}clp{5cm}@{}}
\toprule
\textbf{\#} & \textbf{Lemma} & \textbf{Use} \\
\midrule
1 & \lean{integral\_nonneg} & $\int f \geq 0$ when $f \geq 0$ \\
2 & \lean{integral\_mono\_of\_nonneg} & Bounding integrals \\
3 & \lean{sq\_nonneg} & $x^2 \geq 0$ \\
4 & \lean{Matrix.PosDef.det\_pos} & Gram invertibility \\
5 & \lean{Matrix.nonsing\_inv} & Matrix inversion \\
6 & \lean{Real.log\_le\_sub\_one\_of\_le} & Log inequalities \\
7 & \lean{Real.rpow\_natCast} & Power coercion \\
8 & \lean{Finset.sum\_le\_sum} & Bounding finite sums \\
9 & \lean{riemannZeta\_ne\_zero\_of\_one\_le\_re} & $\zeta(s)\neq 0$ \\
10 & \lean{integral\_univ\_inv\_one\_add\_sq} & $\int(1+t^2)^{-1}=\pi$ \\
11 & \lean{MonotoneBounded.tendsto} & Monotone convergence \\
12 & \lean{Nat.Coprime.gcd\_eq\_one} & GCD properties \\
\bottomrule
\end{tabular}
\end{center}

% ================================================================
\section{Project Structure}
\label{sec:structure}
% ================================================================

\subsection{The Layered Module Pattern}

\begin{enumerate}
  \item \textbf{Layer 0}: Definitions. No theorems, no axioms.
  \item \textbf{Layer 1}: Independent mathematical foundations.
  \item \textbf{Layer 2}: Domain-specific modules (no cross-deps).
  \item \textbf{Layer 3}: Bridges connecting Layer~2 modules.
  \item \textbf{Layer 4}: Assembly (crown theorem).
\end{enumerate}

\subsection{The Axiom Registry Pattern}

Maintain \lean{Axioms.lean} documenting every axiom. Run
\lean{\#print axioms} on the crown theorem after every structural
change. Discrepancies are bugs.

\subsection{Sorry Hygiene}

\begin{enumerate}
  \item Never \lean{sorry} on the critical path.
  \item Document every \lean{sorry} with a comment.
  \item Track sorry count as a project metric.
\end{enumerate}

The Cathedral maintains: 0 sorry on the crown path.

% ================================================================
\section{Tactic Patterns}
\label{sec:tactics}
% ================================================================

\begin{itemize}
  \item \textbf{linarith-norm\_num}: For numerical inequalities.
  \item \textbf{ring-then-bound}: Algebra then analysis.
  \item \textbf{convert-congr}: Escape hatch for type unification.
    Use as last resort.
\end{itemize}

% ================================================================
\section{Deprecated API Traps}
\label{sec:deprecated}
% ================================================================

\begin{center}
\begin{tabular}{@{}lll@{}}
\toprule
\textbf{Deprecated} & \textbf{Replacement} & \textbf{Status} \\
\midrule
\lean{isHermitian\_iff\_isSymmetric} & (none yet) & Kept \\
\lean{Integrable.bdd\_mul'} & \lean{Integrable.bdd\_mul} & Updated \\
\lean{measurable\_of\_continuousOn} & \lean{Continuous.measurable} & Updated \\
\bottomrule
\end{tabular}
\end{center}

When Lean warns ``deprecated,'' check whether the replacement exists
in your Mathlib version.

% ================================================================
\section{Verification Checklist}
\label{sec:checklist}
% ================================================================

\begin{enumerate}
  \item[$\square$] \lean{lake build} passes with zero errors.
  \item[$\square$] \lean{\#print axioms crown\_theorem} shows only
    expected axioms.
  \item[$\square$] No \lean{sorry} on the crown path.
  \item[$\square$] All sorry markers documented.
  \item[$\square$] Ghost axiom detection passes.
  \item[$\square$] Numerical validation agrees with formal claims.
  \item[$\square$] All deprecated API calls documented.
\end{enumerate}

% ================================================================
\section{Conclusion}
% ================================================================

Formalizing analytic number theory in Lean~4 is feasible today, but
requires significant infrastructure work. The main gaps are in
complex analysis (contour integrals, analytic continuation) and
specialized number theory (Mertens function, Dirichlet polynomial
bounds). The Cathedral took 45 days and produced $\sim$60{,}500 lines
across 227 active files.

The foundational analysis shows the proof inhabits ordinary
mathematics ($\ZFC$), requires no large cardinals, and achieves
foundational transparency: every assumption is named, classified,
and tracked. The trust chain is auditable.

\begin{quote}
\emph{A proof is not just a certificate of truth. It is a map of
trust, showing exactly where certainty ends and faith begins.}
\end{quote}

\begin{thebibliography}{99}

\bibitem{lean4}
L.~de~Moura et al., \emph{The Lean~4 Theorem Prover}, CADE-28, 2021.

\bibitem{mathlib}
The Mathlib Community, \emph{Mathlib4},
\url{https://github.com/leanprover-community/mathlib4}.

\bibitem{werner1997}
B.~Werner, \emph{Sets in Types, Types in Sets}, Proc. TACS 1997.

\bibitem{carneiro2019}
M.~Carneiro, \emph{The Type Theory of Lean}, MSc thesis, CMU, 2019.

\bibitem{debruijn1970}
N.G.~de Bruijn, \emph{The Mathematical Language AUTOMATH},
Springer LNCS 125 (1970).

\bibitem{cathedral-math}
J.R.~Gochanour, \emph{The Cathedral: A Machine-Verified Reduction
of the Riemann Hypothesis}, 2026.

\end{thebibliography}

\end{document}
